\subsection{平坦模与挠积}

定理 2. —— 设 E 是一个右 A-模。下列条件是等价的：

\begin{enumerate}
\def\labelenumi{(\roman{enumi})}
\item
  E 是平坦的；
\item
  对于每个左 A-模 F 以及每个整数 n \textgreater{} 0，有
\end{enumerate}

\[
\operatorname{Tor} _ {n} ^ {\mathbf {A}} (\mathbf {E}, \mathbf {F}) = 0;
\]

\begin{enumerate}
\def\labelenumi{(\roman{enumi})}
\setcounter{enumi}{2}
\tightlist
\item
  对于每个有限呈现的单生成左 A-模 F，有
\end{enumerate}

\[
\operatorname{Tor} _ {1} ^ {\mathbf {A}} (\mathbf {E}, \mathbf {F}) = 0;
\]

\begin{enumerate}
\def\labelenumi{(\roman{enumi})}
\setcounter{enumi}{3}
\item
  对于 A 的每个有限生成左理想 a，典范映射 E ⊗\_A a → E 是单射；
\item
  对于如下形式的右 A-模的每个正合序列
\end{enumerate}

\[
0 \rightarrow \mathbf {G} \xrightarrow {v} \mathbf {H} \xrightarrow {w} \mathbf {E} \rightarrow 0,
\]

以及每个左 A-模 F，序列

\[
0 \longrightarrow \mathrm{G} \otimes_ {\mathrm{A}} \mathrm{F} \xrightarrow {v \otimes 1} \mathrm{H} \otimes_ {\mathrm{A}} \mathrm{F} \xrightarrow {w \otimes 1} \mathrm{E} \otimes_ {\mathrm{A}} \mathrm{F} \longrightarrow 0
\]

是正合的。

\begin{enumerate}
\def\labelenumi{(\roman{enumi})}
\item
  \(\Rightarrow\) （ii）：这是 X，第 69 页，命题 5 的推论。
\item
  \(\Rightarrow\) （iii）：这是平凡的。
\item
  \(\Leftrightarrow\) （iv）：每个有限呈现的单生成左 A-模都同构于一个商模 A/a，其中 a 是一个有限生成的左理想，因此由 X，第 72 页，例可知（iii）等价于（iv）。
\item
  \(\Rightarrow\) （i）：由 X，第 8 页，命题 3，X，第 72 页，定理 1，只要 \(\operatorname{Tor}_{1}^{\mathbf{A}}(\mathbf{E},\mathbf{F}) = 0\) 对每个左 A-模 F 成立，E 就是平坦的。若（iii）成立，则当 F 是单生成且有限呈现时，上述条件成立。由 X，第 11 页，命题 7，每个 A-模（相应地，每个单生成 A-模）都是有限呈现模（相应地，有限呈现的单生成模）的滤过归纳极限；因此，由 X，第 70 页，命题 8，可以看出只需证明：若 \(\operatorname{Tor}_{1}^{\mathbf{A}}(\mathbf{E},\mathbf{F}) = 0\) 当 F 是单生成时成立，则当 F 是有限生成时也成立。于是我们对 F 的生成族 \((f_{1},\ldots,f_{n})\) 的基数进行归纳；正合序列
\end{enumerate}

\[
0 \rightarrow \mathrm{A} f _ {1} \rightarrow \mathrm{F} \rightarrow \mathrm{F} / \mathrm{A} f _ {1} \rightarrow 0
\]

给出正合序列

\[
\operatorname{Tor} _ {1} ^ {\mathbf {A}} (\mathrm{E}, \mathrm{A} f _ {1}) \rightarrow \operatorname{Tor} _ {1} ^ {\mathbf {A}} (\mathrm{E}, \mathrm{F}) \rightarrow \operatorname{Tor} _ {1} ^ {\mathbf {A}} (\mathrm{E}, \mathrm{F} / \mathrm{A} f _ {1}),
\]

从而 \(\operatorname{Tor}_{1}^{\mathbf{A}}(\mathbf{E},\mathbf{F})=0\)，因为 \(\operatorname{Tor}_{1}^{\mathbf{A}}(\mathbf{E},\mathbf{A}f_{1})=0\) 且 \(\operatorname{Tor}_{1}^{\mathbf{A}}(\mathbf{E},\mathbf{F}/\mathbf{A}f_{1})=0\)，由归纳假设。

\begin{enumerate}
\def\labelenumi{(\roman{enumi})}
\item
  \(\Rightarrow\) (v)：这是定理1的推论2（X，第72页）。
\item
  \(\Rightarrow\) (iii)：正合序列（X，第50页）
\end{enumerate}

\[
0 \longrightarrow Z _ {0} (\mathrm{E}) \stackrel {i _ {\mathrm{E}}} {\longrightarrow} L _ {0} (\mathrm{E}) \stackrel {p _ {\mathrm{E}}} {\longrightarrow} \mathrm{E} \longrightarrow 0
\]

对每个左A-模F，产生一个正合序列

\[
0 \longrightarrow \operatorname{Tor} _ {1} ^ {\mathrm{A}} (\mathrm{E}, \mathrm{F}) \longrightarrow Z _ {0} (\mathrm{E}) \otimes_ {\mathrm{A}} \mathrm{F} \xrightarrow {i _ {\mathrm{E}} \otimes 1} \mathrm{L} _ {0} (\mathrm{E}) \otimes_ {\mathrm{A}} \mathrm{F} \xrightarrow {p _ {\mathrm{E}} \otimes 1} \mathrm{E} \otimes_ {\mathrm{A}} \mathrm{F} \longrightarrow 0.
\]

若满足(v)，则有 \(\operatorname{Tor}_1^{\mathbf{A}}(\mathbf{E},\mathbf{F}) = 0\)，由此得（iii）。

推论 1. --- 设 \(0 \to E' \to E \to E'' \to 0\) 是右A-模的一个正合序列。假设 \(E''\) 是平坦的。那么E平坦当且仅当 \(E'\) 是平坦的。

设 F 为左 A-模。由于 \(\operatorname{Tor}_i^{\mathbf{A}}(\mathbf{E}^{\prime \prime},\mathbf{F}) = 0\)（\(i = 1,2\)；定理2，(i) \(\Rightarrow\) (ii)），存在一个正合序列

\[
0 \to \mathrm{Tor} _ {1} ^ {\mathbf {A}} (\mathbf {E} ^ {\prime}, \mathbf {F}) \to \mathrm{Tor} _ {1} ^ {\mathbf {A}} (\mathbf {E}, \mathbf {F}) \to 0
\]

由此得出断言（定理2，（i）⇔（iii））。

推论 2. --- 设 \(0 \to E_n \to E_{n-1} \to \cdots \to E_1 \to 0\) 是右A-模的一个正合序列。如果 \(E_i\) 对 \(i = 1, \ldots, n-1\) 是平坦的，则 \(E_n\) 是平坦的。
% semantic-fragments: legacy-input-seam
\relax{}
